% ================================================================
%  sesion20.tex  —  Sesión 20 de 20
%  Exposición y cierre del proyecto
%  Bloque 4: Producto final · Matemáticas 1.º ESO
% ================================================================
\documentclass[aspectratio=169, 12pt, spanish]{beamer}
\usepackage{almeria-beamer}

\renewcommand{\SessionNumber}{20}
\renewcommand{\SessionTitle}{Exposición y cierre del proyecto}
\renewcommand{\SessionName}{Sesión 20}
\renewcommand{\SessionObjective}{%
  Presentar la guía turística matemática ante la clase,
  defender las decisiones matemáticas tomadas, y reflexionar
  sobre los aprendizajes adquiridos a lo largo del proyecto.}
\renewcommand{\BlockName}{Bloque 4: Producto final}

\begin{document}

\begin{frame}[plain]\titlepage\end{frame}

\begin{frame}{Índice}
  \begin{columns}[t]
    \begin{column}{0.5\textwidth}
      \begin{enumerate}
        \item Estructura de la exposición oral
        \item Consejos para presentar matemáticas
        \item Vocabulario matemático requerido
        \item Rúbrica de evaluación oral
        \item Síntesis de los 4 bloques
      \end{enumerate}
    \end{column}
    \begin{column}{0.5\textwidth}
      \begin{enumerate}\setcounter{enumi}{5}
        \item Lo que hemos construido
        \item Preguntas de reflexión final
        \item Ejercicios (evaluar y crear)
        \item Error final: el mito
        \item Cierre del proyecto
      \end{enumerate}
    \end{column}
  \end{columns}
\end{frame}

% ---------------------------------------------------------------
\seccionframe{1. Estructura de la exposición}
% ---------------------------------------------------------------

\begin{frame}{Exposición oral — 5 minutos por equipo}
  \begin{columns}[c]
    \begin{column}{0.55\textwidth}
      \begin{tikzpicture}[
        bloque/.style={
          rounded corners=4pt,
          draw=AlmAzulM, fill=AlmFondo,
          text width=6.5cm, align=left,
          inner sep=7pt, font=\small
        },
        tiempo/.style={
          circle, fill=AlmAmbar, text=AlmAzul,
          font=\bfseries\small, inner sep=3pt,
          minimum size=0.8cm
        }
      ]
        \node[tiempo] (t1) at (0,4.5) {1'};
        \node[bloque, right=0.3cm of t1] {
          \textbf{Introducción}\\
          {\footnotesize Equipo, objetivo del proyecto,
          monumento o zona de Almería elegida como tema central.}
        };
        \node[tiempo] (t2) at (0,3.1) {2'};
        \node[bloque, right=0.3cm of t2] {
          \textbf{2 resultados matemáticos clave}\\
          {\footnotesize Un cálculo completo (área, distancia, etc.)
          y una gráfica explicada. Mostrad vuestros pasos.}
        };
        \node[tiempo] (t3) at (0,1.7) {1'};
        \node[bloque, right=0.3cm of t3] {
          \textbf{Descubrimiento más sorprendente}\\
          {\footnotesize El dato de Almería que más os ha sorprendido
          y por qué tiene importancia matemática.}
        };
        \node[tiempo] (t4) at (0,0.3) {1'};
        \node[bloque, right=0.3cm of t4] {
          \textbf{Reflexión}\\
          {\footnotesize Una cosa que habéis aprendido y una
          que haríais diferente en un segundo intento.}
        };
        % Flecha temporal
        \draw[AlmAmbar, line width=1.5pt, ->]
          (t1.south) -- (t2.north);
        \draw[AlmAmbar, line width=1.5pt, ->]
          (t2.south) -- (t3.north);
        \draw[AlmAmbar, line width=1.5pt, ->]
          (t3.south) -- (t4.north);
      \end{tikzpicture}
    \end{column}
    \begin{column}{0.41\textwidth}
      \begin{notabox}[Materiales permitidos]
        \begin{itemize}
          \item La guía impresa o en pantalla
          \item Fichas de notas (no leer todo)
          \item El mapa de coordenadas para señalar
          \item Las gráficas de la guía en grande
        \end{itemize}
      \end{notabox}
      \vspace{6pt}
      \begin{notabox}[Preguntas del profesor]
        \footnotesize
        Tras la exposición habrá 1--2 preguntas
        que \textbf{cualquier miembro del equipo}
        debe poder responder. No son trampa:
        solo verifican que entendéis lo que
        habéis hecho.
      \end{notabox}
    \end{column}
  \end{columns}
\end{frame}

% ---------------------------------------------------------------
\seccionframe{2. Consejos para presentar}
% ---------------------------------------------------------------

\begin{frame}{Consejos para presentar matemáticas en público}
  \begin{columns}[T]
    \begin{column}{0.48\textwidth}
      \textbf{\color{AlmVerde}\faCheck\ Sí hacer:}
      \begin{itemize}
        \item Mirar a la audiencia, no al papel.
        \item Señalar los números en el mapa o la gráfica
              mientras los explicas.
        \item Usar vocabulario matemático correcto:
              «\textit{calculamos el área con $b \times h$}»
              en vez de «\textit{lo multiplicamos}».
        \item Justificar cada decisión:
              «Elegimos la gráfica lineal porque
              la temperatura es variable continua.»
        \item Hablar despacio en los cálculos:
              la audiencia necesita seguir los pasos.
      \end{itemize}
    \end{column}
    \begin{column}{0.48\textwidth}
      \textbf{\color{AlmTerra}\faTimes\ No hacer:}
      \begin{itemize}
        \item Leer el trabajo en voz alta sin mirarnos.
        \item Decir «este número» sin señalar cuál.
        \item Usar «esto», «aquello», «la cosa»
              en lugar del nombre matemático.
        \item Ignorar las unidades en la presentación oral.
        \item Excusarse continuamente: «Sé que no está bien pero\ldots»\\
              {\footnotesize Si sabes que hay un error,
              corrígelo antes de presentar.}
      \end{itemize}
      \vspace{6pt}
      \begin{notabox}[Ante preguntas difíciles]
        \footnotesize
        «Buena pregunta. Según nuestros datos\ldots»\\
        «No lo hemos calculado, pero podríamos\ldots»\\
        «No estoy seguro/a, lo podría verificar con\ldots»
      \end{notabox}
    \end{column}
  \end{columns}
\end{frame}

% ---------------------------------------------------------------
\seccionframe{3. Vocabulario matemático}
% ---------------------------------------------------------------

\begin{frame}{Vocabulario matemático: ¡úsalo!}
  \begin{columns}[T]
    \begin{column}{0.48\textwidth}
      Debéis usar \textbf{al menos 10} de estos términos
      durante la exposición:
      \vspace{4pt}
      \begin{center}
      \begin{tabular}{ll}
        \toprule
        \textbf{Término} & \textbf{Contexto de uso}\\
        \midrule
        Coordenadas & Localización en el mapa\\
        Perímetro   & Borde de una zona\\
        Área        & Superficie de un espacio\\
        Función     & Relación x→y\\
        Proporcionalidad & k = y/x constante\\
        Variable    & Lo que medimos\\
        Tendencia   & Dirección de una gráfica\\
        Correlación & Relación entre dos variables\\
        Escala      & Conversión mapa--realidad\\
        Magnitud    & Lo que se mide (longitud, etc.)\\
        Perpendicular & Ángulo recto entre calles\\
        Polígono    & Figura cerrada de lados rectos\\
        \bottomrule
      \end{tabular}
      \end{center}
    \end{column}
    \begin{column}{0.48\textwidth}
      \textbf{Ejemplos de frases bien construidas:}\\[6pt]
      \begin{tcolorbox}[colback=AlmVerde!8, colframe=AlmVerde,
                        arc=4pt, boxrule=0.5pt, fontupper=\small]
        «La \textbf{distancia} entre la Alcazaba y el Puerto
        es $d = \sqrt{3^2+3^2} = \sqrt{18} \approx 4{,}2$
        unidades, que con nuestra \textbf{escala}
        de 1\,u = 200\,m equivale a \SI{840}{m}.»
      \end{tcolorbox}
      \vspace{6pt}
      \begin{tcolorbox}[colback=AlmVerde!8, colframe=AlmVerde,
                        arc=4pt, boxrule=0.5pt, fontupper=\small]
        «La gráfica de turistas muestra una
        \textbf{tendencia creciente} de enero a agosto,
        lo que indica una \textbf{correlación positiva}
        con la temperatura.»
      \end{tcolorbox}
    \end{column}
  \end{columns}
\end{frame}

% ---------------------------------------------------------------
\seccionframe{4. Rúbrica de evaluación oral}
% ---------------------------------------------------------------

\begin{frame}{Rúbrica de evaluación de la exposición oral}
  \begin{center}
  \footnotesize
  \begin{tabular}{p{3.5cm} p{2.2cm} p{2.2cm} p{2.2cm} c}
    \toprule
    \textbf{Criterio} & \textbf{Excelente (3)} & \textbf{Bien (2)}
      & \textbf{Mejora (1)} & \textbf{Pts.}\\
    \midrule
    Exactitud matemática\\de los datos
      & Sin errores
      & 1 error menor
      & 2+ errores
      & /3\\[4pt]
    Claridad en la\\explicación de cálculos
      & Paso a paso,\\muy claro
      & Mayormente claro
      & Confuso o\\incompleto
      & /3\\[4pt]
    Vocabulario\\matemático
      & $\geq 10$ términos\\correctos
      & 6--9 términos
      & $<6$ términos
      & /2\\[4pt]
    Calidad visual\\(gráficas, mapas, tablas)
      & Todos completos\\y correctos
      & Algún elemento\\incompleto
      & Varios\\incompletos
      & /2\\
    \midrule
    \multicolumn{4}{r}{\textbf{TOTAL}} & \textbf{/10}\\
    \bottomrule
  \end{tabular}
  \end{center}
  \vspace{4pt}
  \begin{center}
    {\small 9--10: Sobresaliente \quad 7--8: Notable \quad
    5--6: Bien \quad $<5$: Necesita revisión}
  \end{center}
\end{frame}

% ---------------------------------------------------------------
\seccionframe{5. Síntesis: 4 bloques, 20 sesiones}
% ---------------------------------------------------------------

\begin{frame}{Lo que hemos aprendido — Los 4 bloques}
  \begin{columns}[T]
    \begin{column}{0.48\textwidth}
      \begin{tcolorbox}[enhanced, colback=AlmAzulM!12,
        colframe=AlmAzulM, arc=5pt,
        title={\faMapMarkerAlt\enskip Bloque 1: Mapa y coordenadas},
        fonttitle=\bfseries\small\color{AlmBlanco},
        attach boxed title to top left={yshift=-2.2mm, xshift=4mm},
        boxed title style={colback=AlmAzulM, arc=3pt}]
        \footnotesize
        Sesiones 1--4\\
        \textbf{Clave:} localizar cualquier punto con $(x,y)$,
        escala, tipos de información (Ubicación/Medida/Dato)
      \end{tcolorbox}
      \vspace{6pt}
      \begin{tcolorbox}[enhanced, colback=AlmTerra!10,
        colframe=AlmTerra, arc=5pt,
        title={\faCity\enskip Bloque 2: Geometría urbana},
        fonttitle=\bfseries\small\color{AlmBlanco},
        attach boxed title to top left={yshift=-2.2mm, xshift=4mm},
        boxed title style={colback=AlmTerra, arc=3pt}]
        \footnotesize
        Sesiones 5--10\\
        \textbf{Clave:} perímetro, área, polígonos,
        Teorema de Pitágoras para distancias
      \end{tcolorbox}
    \end{column}
    \begin{column}{0.48\textwidth}
      \begin{tcolorbox}[enhanced, colback=AlmVerde!10,
        colframe=AlmVerde, arc=5pt,
        title={\faChartLine\enskip Bloque 3: Datos y gráficas},
        fonttitle=\bfseries\small\color{AlmBlanco},
        attach boxed title to top left={yshift=-2.2mm, xshift=4mm},
        boxed title style={colback=AlmVerde, arc=3pt}]
        \footnotesize
        Sesiones 11--16\\
        \textbf{Clave:} variables, tablas, gráficas,
        proporcionalidad directa $y=kx$
      \end{tcolorbox}
      \vspace{6pt}
      \begin{tcolorbox}[enhanced, colback=BCrear!10,
        colframe=BCrear, arc=5pt,
        title={\faStar\enskip Bloque 4: Producto final},
        fonttitle=\bfseries\small\color{AlmBlanco},
        attach boxed title to top left={yshift=-2.2mm, xshift=4mm},
        boxed title style={colback=BCrear, arc=3pt}]
        \footnotesize
        Sesiones 17--20\\
        \textbf{Clave:} selección, construcción, revisión
        y presentación de la guía turística matemática
      \end{tcolorbox}
    \end{column}
  \end{columns}
\end{frame}

% ---------------------------------------------------------------
\seccionframe{6. Lo que hemos construido}
% ---------------------------------------------------------------

\begin{frame}{¿Qué hemos construido?}
  \begin{columns}[c]
    \begin{column}{0.55\textwidth}
      \begin{resumenbox}
        \textbf{Una guía turística matemática real de Almería:}
        \begin{itemize}
          \item \faMapMarkerAlt\ Mapa con coordenadas cartesianas
                y $\geq 10$ puntos de interés
          \item \faRuler\ Áreas y perímetros de zonas urbanas
          \item \faRoute\ Rutas calculadas con el Teorema de Pitágoras
          \item \faChartBar\ Gráficas de datos reales de Almería
          \item \faTable\ Tablas organizadas con datos verificados
          \item \faEquals\ Proporcionalidades y funciones lineales
          \item \faBookOpen\ Glosario de términos matemáticos
          \item \faLightbulb\ Datos curiosos con significado matemático
        \end{itemize}
      \end{resumenbox}
    \end{column}
    \begin{column}{0.41\textwidth}
      \begin{tikzpicture}[scale=0.7]
        % Guía representada como libro
        \fill[AlmAzul, rounded corners=5pt]
          (0,0) rectangle (5,7);
        \fill[AlmAmbar]
          (0,6.2) rectangle (5,7);
        \node[text=AlmAzul, font=\tiny\bfseries] at (2.5,6.6)
          {ALMERÍA EN CIFRAS Y FORMAS};
        \fill[AlmAzulC!30] (0.3,4.6) rectangle (4.7,6.0);
        \node[text=AlmAzul, font=\tiny, align=center] at (2.5,5.3)
          {mapa + coordenadas};
        \fill[AlmAmbar!20] (0.3,3.5) rectangle (2.2,4.4);
        \node[text=AlmNegro, font=\tiny, align=center] at (1.25,3.95)
          {tabla\\monum.};
        \fill[AlmVerde!20] (2.5,3.5) rectangle (4.7,4.4);
        \node[text=AlmNegro, font=\tiny, align=center] at (3.6,3.95)
          {gráfica\\turistas};
        \foreach \y in {2.7,2.3,1.9,1.5,1.1,0.7}{
          \fill[AlmBlanco!20] (0.3,\y) rectangle (4.7,\y+0.22);
        }
        \node[text=AlmAmbar, font=\tiny] at (2.5,0.3)
          {Matemáticas 1.º ESO · IES Al-Ándalus};
      \end{tikzpicture}
    \end{column}
  \end{columns}
\end{frame}

% ---------------------------------------------------------------
\seccionframe{7. Reflexión final}
% ---------------------------------------------------------------

\begin{frame}{Preguntas para la reflexión final}
  \begin{columns}[T]
    \begin{column}{0.48\textwidth}
      \begin{tcolorbox}[
        enhanced, colback=AlmAmbar!10, colframe=AlmAmbar!70!black,
        title={\faBrain\enskip Reflexiones matemáticas},
        fonttitle=\bfseries\small\color{AlmAzul},
        attach boxed title to top left={yshift=-2.2mm, xshift=4mm},
        boxed title style={colback=AlmAmbar, coltext=AlmAzul, arc=3pt},
        arc=4pt]
        \small
        \begin{enumerate}
          \item ¿Qué herramienta matemática
                (coordenadas, geometría, gráficas)
                ha resultado más \textbf{útil} para
                describir Almería?
          \vspace{4pt}
          \item ¿Qué dato de Almería os ha
                \textbf{sorprendido más} y por qué
                tiene importancia matemática?
          \vspace{4pt}
          \item Si tuvieras que guiar a un turista
                por Almería usando \textbf{solo matemáticas},
                ¿qué tres cosas le mostrarías?
        \end{enumerate}
      \end{tcolorbox}
    \end{column}
    \begin{column}{0.48\textwidth}
      \begin{tcolorbox}[
        enhanced, colback=AlmAzulM!10, colframe=AlmAzulM,
        title={\faStar\enskip Tarjeta de cierre},
        fonttitle=\bfseries\small\color{AlmBlanco},
        attach boxed title to top left={yshift=-2.2mm, xshift=4mm},
        boxed title style={colback=AlmAzulM, arc=3pt},
        arc=4pt]
        \small
        Completa estas frases en tu cuaderno:
        \begin{enumerate}
          \item «La herramienta matemática más útil
                para una guía turística es
                \underline{\hspace{2cm}}
                porque \underline{\hspace{2cm}}.»
          \item «Me sorprendió descubrir que
                \underline{\hspace{3cm}}.»
          \item «La próxima vez mejoraría
                \underline{\hspace{3cm}}.»
          \item «Las matemáticas en una ciudad
                son \underline{\hspace{3cm}}.»
        \end{enumerate}
      \end{tcolorbox}
    \end{column}
  \end{columns}
\end{frame}

% ---------------------------------------------------------------
\seccionframe[BCrear]{8. Ejercicios}
% ---------------------------------------------------------------

\begin{frame}{Ejercicios finales}
  \begin{columns}[T]
    \begin{column}{0.48\textwidth}
      \begin{evaluarbox}
        \footnotesize\dificultad{3}\\[4pt]
        Después de escuchar \textbf{tres exposiciones}
        de otros equipos, evalúa cada una con la rúbrica
        (puntuación numérica + 2 líneas de justificación
        por criterio).\\[6pt]
        ¿Qué aspectos destacan de los mejores trabajos?
        ¿Qué diferencia a una guía «notable» de una
        «sobresaliente»?
      \end{evaluarbox}
    \end{column}
    \begin{column}{0.48\textwidth}
      \begin{crearbox}
        \footnotesize\dificultad{4}\\[4pt]
        \textbf{Proyecto ampliado (opcional):}\\[6pt]
        ¿Y si aplicáramos el mismo método a
        \textbf{otra ciudad}? Elige una ciudad
        (tu ciudad natal, Madrid, Sevilla\ldots)
        y diseña las \textbf{primeras 2 páginas}
        de una guía turística matemática para ella,
        usando exactamente las mismas herramientas:
        mapa con coordenadas, tabla de monumentos
        y una gráfica de datos locales.
      \end{crearbox}
    \end{column}
  \end{columns}
\end{frame}

% ---------------------------------------------------------------
\begin{frame}{Evidencia del día}
  \begin{evidenciabox}
    Entrega final del proyecto — 3 elementos:
    \begin{enumerate}
      \item \textbf{Guía turística matemática} completa
            (versión final revisada con todas las correcciones)
      \item \textbf{Tarjeta de cierre} con las 4 frases completadas
            (reflexión personal)
      \item \textbf{Evaluación de 3 exposiciones} de otros equipos
            con rúbrica completada y justificación escrita
    \end{enumerate}
    \vspace{6pt}
    \begin{center}
      \large\color{AlmAmbar}
      {\Large\faAward}\quad
      \textbf{¡Enhorabuena por completar el proyecto!}
    \end{center}
  \end{evidenciabox}
\end{frame}

% ---------------------------------------------------------------
\begin{frame}{¡Cuidado con este error final!}
  \begin{errorbox}
    \textbf{«Las matemáticas son abstractas y no tienen
    conexión con la vida real»}\\[8pt]
    \incorrecto\ Esta afirmación es el error más grande
    de todos los que hemos visto en 20 sesiones.\\[8pt]
    \correcto\ Hemos demostrado exactamente lo contrario:\\
    Las matemáticas \textbf{sí están} en las calles de Almería,
    en las distancias entre monumentos, en las gráficas
    de turistas, en las áreas de los parques, en la
    proporcionalidad del precio de los transportes y en
    las coordenadas que nos llevan de un lugar a otro.\\[6pt]
    \textit{La ciudad que acabáis de describir matemáticamente
    es la misma ciudad en la que vivís cada día.}
  \end{errorbox}
\end{frame}

% ---------------------------------------------------------------
\begin{frame}{Cierre del proyecto — Almería en Cifras y Formas}
  \begin{center}
    \begin{tikzpicture}
      \fill[AlmAzul, rounded corners=8pt]
        (-5.5,-2.5) rectangle (5.5,2.5);
      \fill[AlmAmbar]
        (-5.5,1.8) rectangle (5.5,2.5);
      \node[text=AlmBlanco, font=\Large\bfseries] at (0,0.6)
        {Almería en Cifras y Formas};
      \node[text=AlmAmbar, font=\normalsize] at (0,-0.1)
        {20 sesiones · 4 bloques · 1 guía turística matemática};
      \node[text=AlmBlanco!70, font=\small] at (0,-0.8)
        {Matemáticas 1.º ESO · \InstituteName};
      \node[text=AlmBlanco!55, font=\footnotesize] at (0,-1.4)
        {\AuthorName · \SchoolYear};
      % Estrellas decorativas
      \foreach \a in {0,60,...,300}{
        \node[text=AlmAmbar, opacity=0.35, font=\scriptsize]
          at (\a:4.2) {$\bigstar$};
      }
    \end{tikzpicture}
  \end{center}
  \vspace{4pt}
  \begin{center}
    {\color{AlmAzulM}\large
    ¡Gracias por vuestro trabajo durante todo el proyecto!}\\[4pt]
    {\color{AlmGris}\footnotesize
    La matemática no es solo un libro de ejercicios.
    La matemática eres tú describiendo el mundo que te rodea.}
  \end{center}
\end{frame}

\end{document}
